\documentclass[a4paper,12pt]{article}
\usepackage{tikz}
\usetikzlibrary{positioning, calc}
% --- Paketi ---
\usepackage[T1]{fontenc}
\usepackage[utf8]{inputenc}
\usepackage[serbian]{babel}
\usepackage{geometry}
\usepackage{amsmath}
\usepackage{xcolor}
\geometry{margin=2.5cm}

% --- Naslovna strana ---
\title{\textbf{Sistem za daljinski nadzor ekrana i zvuka}\\[0.5em]
\large Projekat iz programiranja u Pythonu}
\author{Jovan Savić (2022/0206) \and Vuk Dinić (2022/0689)\\[0.5em]
Elektrotehnički fakultet, Univerzitet u Beogradu}
\date{Oktobar 2025.}

\begin{document}

\maketitle

\thispagestyle{empty} % bez broja strane
\vfill
\begin{center}
Mentor: Prof. dr Milan Bjelica

\newpage
\section*{Sažetak}

U ovom radu prikazan je sistem za daljinski nadzor računara koji omogućava prenos slike ekrana i zvuka u realnom vremenu, kao i osnovnu daljinsku kontrolu ulaznih uređaja. Sistem je zasnovan na arhitekturi tipa klijent–server i sastoji se od dva glavna modula: agenta, koji se pokreće na nadziranom računaru i šalje podatke, i supervisora (servera), koji prima, dekodira i prikazuje video i audio tokove putem grafičkog interfejsa.  

Sistem je koncipiran kao centralizovana mrežna arhitektura u kojoj jedan server (supervisor) može istovremeno da prihvati i nadzire veći broj klijenata (agenata).

\vspace{1cm}
\begin{center}
\textbf{Ključne reči:} supervisor, agent, daljinski nadzor, video streaming, audio prenos, Python, mrežni soketi, TCP
\end{center}
\end{center}
\newpage
\tableofcontents
\thispagestyle{plain}
\newpage

\section{Uvod}

\subsection{Predmet rada}

\textbf{Integrisani sistem za real-time streaming multimedijalnih podataka}
\begin{itemize}
    \item Prenos video i audio signala u realnom vremenu
    \item Sinhronizacija medijskih tokova
    \item Tehnička i praktična analiza performansi
\end{itemize}

\textbf{Tehnička analiza i implementacija sistema za sinhronizovani multimedijalni prenos}
\begin{itemize}
    \item Teorijsko modelovanje prenosa podataka
    \item Implementacija video i audio stream-a
    \item Evaluacija sinhronizacije i kvaliteta signala
\end{itemize}

\textbf{Optimizovani sistem za streaming audio-video sadržaja u realnom vremenu}
\begin{itemize}
    \item Efikasna kompresija i prenos podataka
    \item Sinhronizacija audio i video tokova
    \item Analiza tehničkih performansi i rezultata
\end{itemize}

\subsection{Motivacija i ciljevi}

Svakodnevna upotreba i velika potražnja za striming servisima u edukativne i zabavne svrhe motiviše razvoj pouzdanih sistema za prenos video i audio podataka preko mreže. Proučavanje i implementacija takvih sistema omogućava bolje razumevanje tehnologija iza modernih striming rešenja i poboljšava kvalitet prenosa multimedije.

\textbf{Ciljevi rada:}
\begin{itemize}
    \item Analizirati postojeće metode i protokole za prenos multimedijalnih podataka u realnom vremenu
    \item Implementirati sistem za prenos video i audio signala sa sinhronizacijom medijskih tokova
    \item Izvršiti tehničku evaluaciju performansi sistema kroz teorijsku analizu i praktična merenja
    \item Predložiti optimizacije za poboljšanje kvaliteta prenosa i smanjenje latencije
\end{itemize}

\newpage
\section{Arhitektura sistema}

Sistem je organizovan u centralnu strukturu u kojoj \textbf{supervisor} ima ulogu glavnog kontrolnog čvora. On prima video i audio podatke od svih priključenih \textbf{agenata} (klijenata) u mreži. Svaki agent po priključenju automatski započinje slanje podataka sa svog kraja, čime se omogućava dinamično proširenje mreže bez potrebe za ručnim konfigurisanjem.

Svaki čvor u sistemu, bilo da je u pitanju supervisor ili agent, koristi tri odvojena komunikaciona kanala, odnosno tri mrežna porta:

\begin{itemize}
    \item \textbf{Video port (5000)} – koristi se za prenos video okvira. Podaci se šalju u komprimovanom obliku kako bi se smanjila zauzetost mreže i kašnjenje.
    \item \textbf{Audio port (5001)} – namenjen je prenosu audio signala, sinhronizovanog sa video tokom. Prenos se odvija paralelno sa video kanalom.
    \item \textbf{Kontrolni port (5002)} – služi za razmenu kontrolnih poruka između supervisor-a i agenata, kao što su potvrde o prijemu podataka, sinhronizacija, upravljanje konekcijama i komande za početak ili prekid prenosa.
\end{itemize}

Takva organizacija omogućava da se svaki tip podataka obrađuje nezavisno, čime se povećava stabilnost sistema i pojednostavljuje kontrola mrežnih tokova.


\begin{figure}[h!]
\centering
\begin{tikzpicture}[thick]

% Supervisor u centru
\node[draw, rectangle, fill=gray!15, minimum width=3cm, minimum height=1cm] (supervisor) {Supervisor};

% 4 Agenta raspoređena dalje od supervisora
\node[draw, rectangle, fill=blue!10, minimum width=2.5cm, minimum height=1cm] (a1) at (4,4) {Agent 1};
\node[draw, rectangle, fill=blue!10, minimum width=2.5cm, minimum height=1cm] (a2) at (4,-4) {Agent 2};
\node[draw, rectangle, fill=blue!10, minimum width=2.5cm, minimum height=1cm] (a3) at (-4,4) {Agent 3};
\node[draw, rectangle, fill=blue!10, minimum width=2.5cm, minimum height=1cm] (a4) at (-4,-4) {Agent 4};

% Tri linije po agentu sa malim pomerajem da se ne preklapaju
% Agent 1 (gore desno)
\draw[red, line width=0.8pt] ($(supervisor.north east)+(-0.05,0)$) -- ($(a1.south west)+(-0.05,0)$);
\draw[blue, line width=0.8pt] ($(supervisor.north east)$) -- ($(a1.south west)$);
\draw[green!60!black, line width=0.8pt] ($(supervisor.north east)+(0.05,0)$) -- ($(a1.south west)+(0.05,0)$);

% Agent 2 (dole desno)
\draw[red, line width=0.8pt] ($(supervisor.south east)+(-0.05,0)$) -- ($(a2.north west)+(-0.05,0)$);
\draw[blue, line width=0.8pt] ($(supervisor.south east)$) -- ($(a2.north west)$);
\draw[green!60!black, line width=0.8pt] ($(supervisor.south east)+(0.05,0)$) -- ($(a2.north west)+(0.05,0)$);

% Agent 3 (gore levo)
\draw[red, line width=0.8pt] ($(supervisor.north west)+(-0.05,0)$) -- ($(a3.south east)+(-0.05,0)$);
\draw[blue, line width=0.8pt] ($(supervisor.north west)$) -- ($(a3.south east)$);
\draw[green!60!black, line width=0.8pt] ($(supervisor.north west)+(0.05,0)$) -- ($(a3.south east)+(0.05,0)$);

% Agent 4 (dole levo)
\draw[red, line width=0.8pt] ($(supervisor.south west)+(-0.05,0)$) -- ($(a4.north east)+(-0.05,0)$);
\draw[blue, line width=0.8pt] ($(supervisor.south west)$) -- ($(a4.north east)$);
\draw[green!60!black, line width=0.8pt] ($(supervisor.south west)+(0.05,0)$) -- ($(a4.north east)+(0.05,0)$);

% Legenda - spustiti ispod dijagrama
\node[below=3cm of supervisor, align=center] (legend) {
\begin{tabular}{ll}
\textcolor{red}{\rule{0.5cm}{0.5pt}} & Video kanal (port 5000) \\
\textcolor{blue}{\rule{0.5cm}{0.5pt}} & Audio kanal (port 5001) \\
\textcolor{green!60!black}{\rule{0.5cm}{0.5pt}} & Kontrolni kanal (port 5002)
\end{tabular}
};

\end{tikzpicture}
\caption{Arhitektura sistema: Supervisor i četiri agenta povezani trokanalnom komunikacijom}
\end{figure}

\section{Prenos podataka}

\subsection{Video prenos}

Video prenos se odvija tako što agent periodično hvata sadržaj ekrana i šalje ga \textbf{supervisoru} preko porta \textbf{5000}. Svaki video okvir (frame) se kompresuje korišćenjem \texttt{OpenCV} i \texttt{numpy} biblioteka, čime se značajno smanjuje količina podataka koji se šalju mrežom.  

Prenos se realizuje u sledećim koracima:
\begin{enumerate}
    \item Agent hvata trenutni prikaz ekrana pomoću biblioteke \texttt{mss}.
    \item Dobijena slika se konvertuje u \texttt{JPEG} format i kodira u niz bajtova.
    \item Na početak svakog paketa dodaje se \textbf{header} koji sadrži informaciju o veličini slike u bajtovima.
    \item Supervisor, koji stalno osluškuje port 5000, najpre čita header da bi znao koliko bajtova pripada trenutnom okviru, a zatim preuzima ceo sadržaj slike i dekodira je.
\end{enumerate}

Header je realizovan kao binarni zapis veličine slike u fiksnoj dužini od \textbf{4 bajta}, uz korišćenje Python-ovog modula \texttt{struct} i formata \texttt{!I}, koji omogućava dosledno kodiranje integer vrednosti nezavisno od platforme.  

\textbf{Primer:}
\begin{verbatim}
header = struct.pack("!I", len(frame_data))
\end{verbatim}

Korišćenje headera omogućava da se tačno zna gde jedan video okvir završava, što je ključno za ispravno dekodiranje slike na strani prijemnika i sprečava gubitak sinhronizacije između paketa.  

Tipična veličina prenete slike iznosi između \textbf{30 kB i 150 kB}, u zavisnosti od rezolucije (npr. 640x480 piksela) i nivoa kompresije. Ovakav pristup omogućava stabilan i kontinuiran video prenos sa malim kašnjenjem, čak i u uslovima ograničenog protoka mreže.

\subsection{Audio prenos}

Audio prenos se odvija tako što agent hvata zvuk sa svog mikrofona i šalje ga \textbf{supervisoru} preko porta \textbf{5001}. Zvuk se obrađuje u malim blokovima (frames) i kompresuje se kako bi zauzeo manje mrežne propusnosti i omogućio real-time prenos.

Prenos se realizuje u sledećim koracima:
\begin{enumerate}
    \item Agent prikuplja audio signal pomoću biblioteke \texttt{pyaudio}, u blokovima od npr. 1024 uzorka.
    \item Svaki blok audio podataka se kodira u binarni niz bajtova PCM tehnikom.
    \item Na početak svakog bloka dodaje se \textbf{header} koji sadrži dužinu bloka u bajtovima.
    \item Supervisor, koji osluškuje port 5001, najpre čita header da bi znao koliko bajtova pripada trenutnom audio bloku, a zatim preuzima podatke i dekodira ih za reprodukciju ili dalje procesiranje.
\end{enumerate}

Header je realizovan kao fiksna 4-bajtna vrednost, koristeći Python modul \texttt{struct} i format \texttt{!I}, što omogućava dosledno kodiranje veličine bloka nezavisno od platforme:

\begin{verbatim}
header = struct.pack("!I", len(audio_data))
\end{verbatim}

Korišćenje headera omogućava supervisoru da pravilno prepozna granice audio blokova i sinhronizuje reprodukciju sa video tokom. Tipična veličina audio bloka zavisi od uzorka po sekundi (npr. 44100 Hz) i dužine bloka (1024 ili 2048 uzoraka), a prenosi se kontinuirano kako bi se obezbedio real-time prenos zvuka sa minimalnom latencijom.
\begin{figure}[h!]
\centering
\begin{tikzpicture}[scale=1.2]

% Os x i y
\draw[->] (0,0) -- (8,0) node[right] {Vreme};
\draw[->] (0,-2.5) -- (0,2.5) node[above] {Napon};

% Sinusoidalni signal
\draw[thick, domain=0:8, smooth, variable=\x, blue] 
    plot ({\x}, {2*sin(180*\x/3.14)});

% Sampling points i PCM linije
\foreach \x in {0.5,1.5,2.5,3.5,4.5,5.5,6.5,7.5} {
    \pgfmathsetmacro{\y}{2*sin(180*\x/3.14)};
    \draw[red, thick] (\x,0) -- (\x,\y);
    \fill[red] (\x,\y) circle (0.05);
}

% Ograničen broj kvantizacionih nivoa za ilustraciju
\foreach \i in {-1.75,-1.25,-0.75,-0.25,0.25,0.75,1.25,1.75} {
    \draw[black!70, dashed] (0,\i) -- (8,\i);
}

\end{tikzpicture}

% Legenda ispod slike
\begin{center}
\begin{tabular}{ll}
\textcolor{blue}{\rule{0.5cm}{0.5pt}} & Sinus signal \\
\textcolor{red}{\rule{0.5cm}{0.5pt}} & PCM sampling \\
\textcolor{black!70}{\rule{0.5cm}{0.5pt}} & Kvantizacioni nivoi (primer)
\end{tabular}
\end{center}

\caption{Princip PCM: sinusoidalni signal, uzorci i nekoliko kvantizacionih nivoa za ilustraciju}
\end{figure}

\subsubsection*{Objašnjenje PCM-a}

Pulse Code Modulation (PCM) funkcioniše na sledeći način:  
\begin{itemize}
    \item Originalni analogni signal se u određenim vremenskim intervalima uzorkuje.  
    \item Svaka uzeta vrednost u periodi se \textbf{zaokružuje na najbliži kvantizacioni nivo}.  
    \item Nakon zaokruživanja, vrednost se \textbf{binarno kodira} u fiksni broj bita koji predstavlja nivo.  
    \item Na taj način se analogni signal pretvara u digitalni tok podataka koji se može prenositi i obrađivati.
\end{itemize}
\subsection{Sinhronizacija}

U sistemu je implementirana sinhronizacija između audio, video i kontrolnih tokova koristeći \textit{odvojene niti} za svaki kanal u okviru svakog čvora (supervisor i agent).  

Pristup zajedničkim resursima (poput bafera i socket-a) je \textit{zaštićen mehanizmima međusobnog isključivanja (mutex/lock)} kako bi se izbegli konflikti i gubici podataka.  

Za sinhronizaciju i obradu podataka korišćeni su standardni algoritmi za \textit{redove poruka (queue), blokiranje niti (thread blocking) i semafore (semaphores)}, što omogućava precizan i stabilan prenos podataka.  

Ova arhitektura omogućava:
\begin{itemize}
    \item Paralelni prenos audio i video podataka bez međusobnog zastoja.  
    \item Održavanje real-time performansi uz minimalno kašnjenje.  
    \item Jednostavno skaliranje sistema sa više agenata.
\end{itemize}

\newpage
\section{Računica i tehnička analiza}

Za početne i referentne vrednosti korišćene u sistemu uzimane su što bliže industrijskim standardima kako bi performanse bile relevantne i praktične.

\subsection{Audio prenos}

Trenutni parametri audio prenosa su:
\begin{itemize}
    \item RATE = 22050 Hz
    \item CHANNELS = 1
    \item CHUNK = 2048
    \item Format = 16-bit PCM
\end{itemize}

Bitrate po kanalu se računa kao:
\[
\text{Bitrate} = \text{RATE} \times \text{CHANNELS} \times \text{Bajt po uzorku} \times 8
\]
\[
\text{Bitrate} = 22050 \times 1 \times 2 \times 8 = 352800~\text{bit/s} \approx 343~\text{kb/s}
\]

Trajanje jednog audio bloka (CHUNK):
\[
t_\text{chunk} = \frac{\text{CHUNK}}{\text{RATE}} = \frac{2048}{22050} \approx 0.093~\text{s}
\]

Dakle, sistem šalje oko 10 audio blokova u sekundi.

\subsubsection*{Proširenje na 44.1 kHz i stereo}

Za stereo audio i standardnu frekvenciju od 44.1~kHz:
\[
\text{Bitrate} = 44100 \times 2 \times 2 \times 8 = 1\,411\,200~\text{bit/s} \approx 1.34~\text{Mb/s}
\]

Povećanje je linearno sa brojem kanala.  
Sistem se može proširiti tako da:
\begin{itemize}
    \item svaki audio kanal koristi zaseban TCP socket
    \item prenos se odvija paralelno i sinhronizovano
    \item omogućava stereo i višekanalni (surround) prenos bez gubitka kvaliteta
\end{itemize}

Takvo rešenje može dostići standard \textbf{44.1 kHz stereo} kvalitet zvuka uz minimalno povećanje latencije.

\subsection{Video prenos - nekompresovani RGB}

Trenutni parametri video prenosa u projektu:

\begin{itemize}
    \item Rezolucija: $640 \times 480$ piksela
    \item Broj frejmova: 30 fps
    \item Boje po pikselu: 24 bita (RGB)
    \item Kompresija: trenutno nije korišćena (nekompresovan prenos)
\end{itemize}

\subsubsection*{Bitrate nekompresovanog video signala}

Bitrate se računa kao:
\[
\text{Bitrate} = \text{Width} \times \text{Height} \times \text{Bojanje po pikselu} \times \text{FPS}
\]

Za naše parametre:
\[
\text{Bitrate} = 640 \times 480 \times 24 \times 30 = 221\,184\,000~\text{bit/s} \approx 221~\text{Mbps}
\]

\subsubsection*{Trajanje jednog frejma}

\[
t_\text{frame} = \frac{1}{\text{FPS}} = \frac{1}{30} \approx 0.033~\text{s}
\]

Dakle, sistem šalje oko 30 frejmova u sekundi.

\subsubsection*{Ukupna propusnost}

\[
\text{Ukupna propusnost video + audio} \approx 221~\text{Mbps} + 0.34~\text{Mbps} \approx 221.34~\text{Mbps}
\]

\section{Performanse}

Sva merenja su vršena na računarima u prostorijama fakulteta, koristeći lokalnu mrežu sa jednim supervisorom i pet agenata. Tokom testiranja izmereni su sledeći pokazatelji performansi:

\begin{table}[htb]
\centering
\caption{Performanse sistema sa 5 agenata}
\medskip
\label{tab:performanse}
\begin{tabular}{l | c}
\hline
Parametar & Vrednost \\ 
\hline
\hline
Latency (video) & 30-50 ms \\ 
Latency (audio) & 90-100 ms \\ 
Propusnost (video) & 221 Mbps \\ 
Propusnost (audio) & 343 kbps \\ 
CPU utilizacija (supervisor) & 40-50\% \\ 
CPU utilizacija (agent) & 15-25\% \\ 
Memorijsko opterećenje (supervisor) & 150-200 MB \\ 
Memorijsko opterećenje (agent) & 50-80 MB \\ 
\hline
\end{tabular}
\end{table}

\noindent
\noindent
\textbf{Način reprodukcije rezultata:}

\begin{itemize}
    \item \textbf{Latencija video signala:} Može se meriti pokretanjem stoperice na računaru supervisora i istovremenim beleženjem ekrana agenta preko screen share-a. Razlika između prikazanog trenutka i stvarnog prenosa omogućava određivanje kašnjenja slike.
    
    \item \textbf{Latencija audio signala:} Precizno se može odrediti tako što se audio signal sa agenta i sa supervisora snimi, a zatim u programu poput \textit{Audacity} uporede karakteristični pikovi zvuka. Razlika u vremenu između tih pikova predstavlja latenciju audio prenosa.
    
    \item \textbf{CPU opterećenje:} Očitava se direktno iz Task Manager-a tokom normalnog rada sistema, prateći kako procesor reaguje na paralelni prenos podataka od svih pet agenata.
    
    \item \textbf{Memorijsko opterećenje:} Procena količine memorije zauzete za baferovanje audio i video podataka, što pomaže u proveri da li je sistem sposoban da održava real-time prenos.
    
    \item \textbf{Propusnost:} Izračunava se prema veličini paketa i broju frejmova/audio blokova poslatih u sekundi, omogućavajući kontrolu ukupnog protoka podataka kroz mrežu.
\end{itemize}

\section{Rezultati i diskusija}

Analiza eksperimentalnih rezultata ukazuje da sistem funkcioniše stabilno u realnom vremenu i da je arhitektura sa odvojenim nitima za audio, video i kontrolne tokove efikasna. Zabeležena latencija i opterećenje resursa pružaju osnovu za identifikaciju potencijalnih optimizacija.

Jedno od ključnih poboljšanja je uvođenje **video kompresije** (npr. H.264), što bi omogućilo smanjenje propusnosti, povećanje broja frejmova ili dodavanje novih agenata bez značajnog povećanja mrežnog opterećenja. Time bi se poboljšala skalabilnost sistema i smanjila upotreba CPU resursa.

Takođe, prelazak sa **TCP na UDP protokol** za prenos audio i video podataka mogao bi dodatno smanjiti latenciju, jer UDP ne koristi potvrde i retransmisije. To bi omogućilo brži prenos u realnom vremenu, ali bi zahtevalo dodatne mehanizme za korekciju gubitaka paketa i očuvanje sinhronizacije tokova.

Dodatno, rezultati ukazuju na mogućnost poboljšanja sinhronizacije i balansiranja resursa između više agenata, što bi doprinelo stabilnijem radu u mrežama sa većim brojem korisnika. Ova analiza pruža jasne smernice za buduća unapređenja sistema u cilju veće efikasnosti i skalabilnosti.
\newpage
\section{Pokretanje i upravljanje sistemom}


\subsection{Pokretanje projekta}

Sve instrukcije za pokretanje sistema preko komandne linije, uključujući komande za pokretanje glavnog skripta, instalaciju zavisnosti i kompajliranje fajlova, detaljno su opisane u \texttt{README} fajlu. Za praktičan primer, sledeće komande mogu biti korišćene:

\begin{center}
\begin{tabular}{|p{14cm}|}  % povećana širina kolone
\hline
\begin{verbatim}
# Instalacija zavisnosti
pip install -r requirements.txt

# Pokretanje glavnog skripta
python main.py

# Alternativno pokretanje preko skripte (Windows)
run_supervisor.bat

# Alternativno pokretanje preko skripte (Linux)
./run_supervisor.sh

# Pokretanje build-portable skripte za kreiranje exe fajlova
./build-portable
\end{verbatim} \\
\hline
\end{tabular}
\end{center}

Skripta \texttt{build-portable} omogućava:
\begin{itemize}
    \item Instalaciju svih zavisnosti.
    \item Kompajliranje Python skripti.
    \item Kreiranje izvršnih fajlova (\texttt{.exe}) za \textit{Agent} i \textit{Supervisor}, koji su spremni za pokretanje na bilo kojoj mašini.
\end{itemize}

Korisnik može birati između direktnog pokretanja glavnog skripta, korišćenja priloženih skripti ili potpuno automatskog build procesa, što čini proces pokretanja jednostavnim i reproduktivnim.

\end{document}
